\documentclass[12pt]{article}


\usepackage[dvips,letterpaper,margin=0.75in,bottom=0.5in]{geometry}
\usepackage{cite}
\usepackage{slashed}
\usepackage{graphicx}
\usepackage{amsmath}
\usepackage{braket}
\usepackage{feynman}

\DeclareMathOperator{\Tr}{Tr}
\newcommand{\gv} {\ensuremath{g_{\mathrm v}}}
\newcommand{\ga} {\ensuremath{g_{\mathrm a}}}
\newcommand{\Ps} {\ensuremath{\slashed{P}}}
\newcommand{\Qs} {\ensuremath{\slashed{Q}}}
\newcommand{\Rs} {\ensuremath{\slashed{R}}}
\newcommand{\Ss} {\ensuremath{\slashed{S}}}
\newcommand{\ps}{\ensuremath{\slashed{p}}}
\newcommand{\GeV} {\ensuremath{\mathrm{Ge\kern -0.08em V}}}

\begin{document}

\title{Take-Home Problem for Final Exam}

\maketitle

\noindent

This assignment is open book and open notes.  The TAs may not provide
any help.  You will be graded subjectively on the strength, accuracy,
and readibility (including neatness) of your argument.  Aim for a
terse but complete exposition, and use the techniques and conventions
as developed in the course.  To discourage throwing everything at the
wall and hoping for partial credit, I will deduct points for pointless
diversions.

{\bf Warning:} do not attempt to use internet solutions for this
problem.  It was choosen in part because several solutions to similar
problems online have subtle mistakes that do not impact the final
result, and which I will be alert for.  ChatGPT's attempt at it was
simply adorable.  If you are caught cheating in this course, you will
fail the entire course.



\begin{figure}[htbp]
\begin{center}
\begin{feynman}
    \fermion{1.00, 1.00}{0.0, 0.2}
    \fermion{0.0, 1.8}{1.00, 1.00}
    \electroweak{1.00, 1.00}{2.00, 1.00}
    \fermion{3.00, 0.2}{2.00, 1.00}
    \fermion{2.00, 1.00}{3.00, 1.8}
    \node at (0,0) {$\bar{v}(b)$};
    \node at (0,2) {$u(a)$};
    \node at (3,2) {$\bar{u}(1)$};
    \node at (3,0) {$v(2)$};
    \node at (1.5,1.2) {$\gamma$};
    \node at (0,0.4) {$e$};
    \node at (0,1.6) {$e$};
    \node at (3.0,0.4) {$\mu$};
    \node at (3.0,1.6) {$\mu$};
\end{feynman}
\caption{\label{fig:feyn} Leading order Feynman diagram for the process $e^- + e^+ \to \gamma \to \mu^- + \mu^+$.  The particle plane-wave bi-spinor is $u$ and the anti-particle plane-wave bi-spinor is $v$.}
\end{center}
\end{figure}
~\\[5pt]

We will consider the process $e^- + e^+ \to \mu^- + \mu^+$ at leading
order, as shown in Fig.~\ref{fig:feyn}.  In Homework 7 you found the
amplitude for this process to be:
\begin{eqnarray*}
\mathcal{M}  &=&  -\frac{g_e^2}{(p_a+p_b)^2} \; \overline{v}(b) \, \gamma^\mu
  \, u(a)  \; \bar{u}(1) \, \gamma_\mu \, v(2) \; 
\end{eqnarray*}
where:
$$\gamma_\mu \equiv g_{\mu\nu} \gamma^\nu$$
In this problem, we will consider an experiment that is capable of
supplying polarized beams, so that the incident electrons and positrons
can be prepared with a known helicity.  We will consider the limit
that $m_e,m_\mu \ll E$.  In this limit the helicity eigenstates are
identical to the chirality eigenstates, which for particles are:
$$\def\arraystretch{1.7}
u_R = N \begin{pmatrix} C  \\ S e^{i\phi} \\ C \\ S e^{i\phi} \end{pmatrix}, \hspace{1cm}
u_L = N \begin{pmatrix} -S \\ C e^{i\phi} \\ S \\ -C e^{i\phi} \end{pmatrix},
$$
and for anti-particles are:
$$\def\arraystretch{1.7}
v_R = N \begin{pmatrix} S \\ -C e^{i\phi} \\ -S \\ C e^{i\phi} \end{pmatrix}, \hspace{1cm}
v_L = N \begin{pmatrix} C \\  S e^{i\phi} \\  C \\ S e^{i\phi} \end{pmatrix},
$$
with
$$C \equiv \cos \left( \frac{\theta}{2} \right), \hspace{1cm } S \equiv \sin \left( \frac{\theta}{2} \right)$$
and $N = \sqrt{E}$.\\[5pt]

\noindent
{\bf Part 1:}  For particle (a) take $\theta=0$ and $\phi=0$ and for anti-particle (b) take $\theta=\pi$ and $\phi=\pi$.  Work out the chirality eigenstates for the particle (a) as:
$$\def\arraystretch{1.7}
u_R(a) = \sqrt{E} \begin{pmatrix} 1 \\ 0 \\ 1 \\ 0 \end{pmatrix}, \hspace{1cm}
u_L(a) = \sqrt{E} \begin{pmatrix} 0 \\ 1 \\ 0 \\ -1 \end{pmatrix},
$$
and for the anti-particle (b) as:
$$\def\arraystretch{1.7}
v_R(a) = \sqrt{E} \begin{pmatrix} 1 \\ 0 \\ -1 \\ 0 \end{pmatrix}, \hspace{1cm}
v_L(a) = \sqrt{E} \begin{pmatrix} 0 \\ -1 \\ 0 \\ -1 \end{pmatrix},
$$
~\\[5pt]

\newpage

\noindent
{\bf Part 2:}  For particle (1) take $\phi=0$ and work out the chirality eigenstates as:
$$\def\arraystretch{1.7}
u_R(1) = \sqrt{E} \begin{pmatrix}  C \\ S \\ C \\  S \end{pmatrix}, \hspace{1cm}
u_L(1) = \sqrt{E} \begin{pmatrix} -S \\ C \\ S \\ -C \end{pmatrix}
$$
~\\[5pt]

\noindent
{\bf Part 3:}  For anti-particle (2) take $\phi=\pi$ and replace $\theta \to \pi - \theta$ to work out the chirality eigenstates:
$$\def\arraystretch{1.7}
v_R(2) = \sqrt{E} \begin{pmatrix}  C \\  S \\ -C \\  -S \end{pmatrix}, \hspace{1cm}
v_L(2) = \sqrt{E} \begin{pmatrix}  S \\ -C \\  S \\  -C \end{pmatrix}
$$
~\\[5pt]

In the amplitude, we encounter vector currents of the form:
\begin{eqnarray*}
  J^\mu \; &=& \; \bar{v}_L(b) \; \gamma^\mu \; u_R(a) \\
  P^\mu \; &=& \; \bar{u}_R(1) \; \gamma^\mu \; v_L(2) \\
  Q^\mu \; &=& \; \bar{u}_L(1) \; \gamma^\mu \; v_R(2) \\
\end{eqnarray*}
\noindent
{\bf Part 4:} Calculate the currents $J^\mu$, $P^\mu$ and $Q^\mu$.  Partial answer:
\begin{eqnarray*}
  (J^0,J^1,J^2,J^3) \; &=& \; 2E(0,-1,-i,0) \\
  (P^0,P^1,P^2,P^3) \; &=& \; 2E(0,-\cos\theta,i,\sin\theta) \\
  (Q^0,Q^1,Q^2,Q^3) \; &=& \; ??? \\
\end{eqnarray*}
(To be clear:  You must still provide the calculation even for the currents with an answer provided!)\\[5pt]    

\noindent
{\bf Part 5:} Consider the process with the specific chiralities $e^-(R) + e^+(L) \to \mu^-(R) + \mu^+(L)$ which has amplitude:
$$\mathcal{M}_{RL \to RL} = -\frac{g_e^2}{(p_a+p_b)^2} \; J^\mu P^\nu g_{\mu\nu}$$
Show that:
$$\mathcal{M}_{RL \to RL} = g_e^2 (1 + \cos \theta)$$

\noindent
{\bf Part 6:} 
Consider the process with the specific chiralities $e^-(R) + e^+(L) \to \mu^-(L) + \mu^+(R)$ which has amplitude:
$$\mathcal{M}_{RL \to LR} = -\frac{g_e^2}{(p_a+p_b)^2} \; J^\mu Q^\nu g_{\mu\nu}$$
Calculate $\mathcal{M}_{RL \to LR}$ as in the previous part.\\[5pt]

\noindent
{\bf Part 7:} 
Calculate the amplitude for the process $e^-(R) + e^+(L) \to \mu^- + \mu^+$, where the spins of the outgoing particles are not specified.  In lecture, we showed that the only nonzero contributions will be:
$$|\mathcal{M}_{RL \to XX}|^2 = |\mathcal{M}_{RL \to RL}|^2 + |\mathcal{M}_{RL \to LR}|^2$$
Compare your result to what we obtained in Homework 7:
$$\braket{|\mathcal{M}|^2} = g_e^4 (1 + \cos^2 \theta)$$
Explain why your result makes sense physically.

\end{document}

